% =====================================================================
\chapter{El coste de la inferencia: \emph{prefill}, decodificación y caché KV}
\label{tema:2}
% =====================================================================

Generar una respuesta tiene dos fases con perfiles de coste opuestos: el \emph{prefill}, que procesa el prompt completo en paralelo, y la \emph{decodificación}, que produce los tokens uno a uno. Este tema demuestra que la caché de claves y valores (caché KV) evita recalcular el pasado sin alterar la salida, calcula su tamaño y aplica el modelo \emph{roofline} del módulo~17 para explicar por qué la decodificación está limitada por el ancho de banda de memoria y cómo el tamaño del lote intercambia latencia por rendimiento. Se siguen la clase~10 de CS336~\cite{cs336}, la sección~2 de Pope \emph{et al.}~\cite{pope}, la sección~3 de Kwon \emph{et al.}~\cite{kwon} y el artículo de Shazeer sobre la atención multiconsulta~\cite{shazeer}.

\section{Las dos fases de la generación}

\begin{definicion}[\emph{Prefill} y decodificación]\label{def:fases}
Sea un prompt $x_{1:n}$. La fase de \emph{prefill} ejecuta una pasada del Transformer sobre las $n$ posiciones a la vez, como en el entrenamiento, y obtiene la distribución del primer token nuevo. La fase de \emph{decodificación} ejecuta, para cada token nuevo, una pasada sobre una sola posición, la del último token emitido.
\end{definicion}

Por el módulo~16, una pasada hacia delante cuesta unas $2N$ operaciones por token y posición, donde $N$ es el número de parámetros (sin contar la atención sobre el contexto, que se trata en el \cref{prop:intensidad-capas}).

\addtocounter{definicion}{1}%

\section{La caché KV}

En un Transformer causal (módulos~13 y~14), sean $h^{(\ell)}_i$ el estado de la posición $i$ a la salida de la capa $\ell$, con $h^{(0)}_i$ la incrustación del token $x_i$ y su posición. Cada capa calcula, para cada posición, una consulta, una clave y un valor por cabeza, $q_i=h^{(\ell-1)}_iW_Q$, $k_i=h^{(\ell-1)}_iW_K$, $v_i=h^{(\ell-1)}_iW_V$, y la salida de la atención en $t$ es
\[
a_t=\softmax\Bigl(\frac{q_t\,[k_1;\dots;k_t]\T}{\sqrt{\dhd}}\Bigr)[v_1;\dots;v_t],
\]
seguida de las conexiones residuales, la normalización y el perceptrón multicapa, que actúan posición a posición.

\begin{lema}[Causalidad de los estados]\label{lem:causal}
Para toda capa $\ell$ y toda posición $i$, el estado $h^{(\ell)}_i$ depende solo de $x_{1:i}$.
\end{lema}

\begin{teorema}[Corrección de la caché KV]\label{teo:cache}
Sea un procedimiento que, en cada paso de decodificación, calcula $q$, $k$ y $v$ únicamente para la posición nueva $t$ en cada capa, añade $(k_t,v_t)$ a una caché que contiene los pares $(k_j,v_j)$, $j<t$, calculados en pasos anteriores, y atiende sobre toda la caché. Entonces sus salidas coinciden exactamente (en aritmética exacta) con las de una pasada completa sobre $x_{1:t}$.
\end{teorema}

En coma flotante pueden aparecer diferencias de redondeo, porque el orden de las operaciones cambia; la distribución que define el modelo es la misma.

\begin{proposicion}[Pasadas de posición con y sin caché]\label{prop:pasadas}
Sea un prompt de $n$ tokens del que se generan $T$ tokens. Sin caché, el número total de posiciones procesadas por las capas no atencionales es $Tn+T(T-1)/2$; con caché es $n+T-1$.
\end{proposicion}

\addtocounter{figure}{1}%

\begin{proposicion}[Tamaño de la caché KV]\label{prop:tam-cache}
Sea un modelo con $L$ capas, $\nkv$ cabezas de clave y valor de dimensión $\dhd$, y números de $s$ bytes. La caché ocupa
\[
c_{\mathrm{kv}}=2\,L\,\nkv\,\dhd\,s\quad\text{bytes por token,}
\]
y una secuencia de $S$ tokens ocupa $S\,c_{\mathrm{kv}}$ bytes.
\end{proposicion}

\addtocounter{definicion}{1}%

\begin{observacion}[Atención multiconsulta y agrupada]
La caché es proporcional a $\nkv$. La atención multiconsulta (MQA)~\cite{shazeer} usa $\nkv=1$ y la agrupada (GQA, módulo~18) un valor intermedio; respecto a la atención multicabeza con $n_h$ cabezas, la caché se divide por $n_h/\nkv$. Pope \emph{et al.} muestran que la atención multiconsulta permite contextos hasta 32 veces más largos con la misma memoria en PaLM 540B~\cite{pope}.
\end{observacion}

\section{Intensidad aritmética de la inferencia}

Se recuerda del módulo~17 que la intensidad aritmética de una operación es $I=F/Q$ (FLOP por byte transferido) y que, en un dispositivo con rendimiento $\Ppico$ y ancho de banda $\Bmem$, la operación está limitada por la memoria si $I<\Icrit=\Ppico/\Bmem$.

\begin{proposicion}[Intensidad de las capas]\label{prop:intensidad-capas}
Sea un lote de $B$ secuencias en bf16 ($s=2$) del que se procesan $T$ posiciones nuevas por secuencia, con $S$ posiciones ya presentes en la caché.
\begin{enumerate}[label=(\alph*)]
  \item En un perceptrón con tres matrices de $D\times F$ (como el SwiGLU), $F_{\mathrm{op}}=6BTDF$ y $Q=4BTD+4BTF+6DF$. Si $BT\ll D,F$, la intensidad es $I\approx BT$.
  \item En la atención, con dimensión total $D$ de consultas, claves y valores, $F_{\mathrm{op}}=4BSTD$ y $Q=4BSD+4BTD$, de modo que $I=ST/(S+T)$, independiente de $B$. En el \emph{prefill} ($T=S$), $I=S/2$; en la decodificación ($T=1$), $I=S/(S+1)<1$.
\end{enumerate}
\end{proposicion}

La diferencia entre ambas capas es estructural: en el perceptrón todas las secuencias del lote comparten los mismos pesos, que se leen una vez; en la atención cada secuencia tiene su propia caché, de modo que agrupar secuencias no reutiliza datos~\cite{cs336}.

\begin{teorema}[Cota de memoria de la decodificación]\label{teo:cota-memoria}
Sea un modelo de $N$ parámetros de $s$ bytes y sea un paso de decodificación con un lote de $B$ secuencias de $S$ tokens de contexto cada una. Si los pesos y la caché se leen al menos una vez por paso, el tiempo del paso cumple
\[
t_{\mathrm{paso}}\ \ge\ \frac{sN+B\,S\,c_{\mathrm{kv}}}{\Bmem}.
\]
En particular, con $B=1$ y caché despreciable, la velocidad de generación cumple $\text{tokens/s}\le\Bmem/(sN)$.
\end{teorema}

\addtocounter{definicion}{1}%

\section{El lote: latencia frente a rendimiento}

\begin{proposicion}[Compromiso del lote]\label{prop:lote}
En las condiciones del \cref{teo:cota-memoria}, supóngase que el tiempo de un paso es exactamente $t(B)=t_0+Bt_1$, con $t_0=sN/\Bmem$ y $t_1=S\,c_{\mathrm{kv}}/\Bmem$. Entonces la velocidad por usuario $1/t(B)$ decrece con $B$, el rendimiento total $R(B)=B/t(B)$ crece con $B$ y
\[
\lim_{B\to\infty}R(B)=\frac{1}{t_1}=\frac{\Bmem}{S\,c_{\mathrm{kv}}}.
\]
\end{proposicion}

El lote amortiza la lectura de los pesos, pero no la de la caché. La \cref{tab:lote} lo ilustra con Llama~2 7B en fp16 ($14$~GB de pesos), secuencias de $2000$ tokens ($1{,}05$~GB de caché cada una) y $\Bmem=2$~TB/s; el límite del rendimiento es $\Bmem/(S\,c_{\mathrm{kv}})\approx 1907$ tokens/s. Además, la memoria total limita $B$: con 64 secuencias la caché ya ocupa $67$~GB.

\begin{table}[h]
\centering\small
\begin{tabular}{@{}rrrrr@{}}
\toprule
Lote $B$ & Caché total (GB) & Tiempo por paso (ms) & Tokens/s totales & Tokens/s por usuario \\
\midrule
1 & 1{,}05 & 7{,}52 & 133 & 133 \\
8 & 8{,}39 & 11{,}19 & 715 & 89{,}3 \\
32 & 33{,}55 & 23{,}78 & 1346 & 42{,}1 \\
64 & 67{,}11 & 40{,}55 & 1578 & 24{,}7 \\
\bottomrule
\end{tabular}
\caption{Cota de memoria del \cref{teo:cota-memoria} para Llama~2 7B en fp16 con contexto de $2000$ tokens y $\Bmem=2$~TB/s.}
\label{tab:lote}
\end{table}

\begin{observacion}[Cuándo deja de mandar la memoria]
Respecto a los pesos, la intensidad de un paso con lote $B$ es $\approx B$ FLOP/byte en bf16 (\cref{prop:intensidad-capas}), de modo que haría falta $B\approx\Icrit$ para quedar limitado por el cómputo; la caché, con intensidad menor que~1, lo impide antes. Pope \emph{et al.} cuantifican este compromiso en PaLM 540B: el \emph{prefill} de lotes grandes alcanza un $76\,\%$ de utilización de FLOP, mientras que la decodificación con lote pequeño logra una latencia de 29~ms por token con pesos en int8~\cite{pope}. Todo lo que reduzca bytes acelera la decodificación: menos cabezas de clave y valor, cuantización (\cref{tema:4,tema:5}) y más tokens por lectura de pesos (\cref{tema:3}).
\end{observacion}

